\documentclass[11pt]{amsart}
\usepackage[T1]{fontenc}
\usepackage{lmodern}
\usepackage{amsmath,amssymb,amsthm,mathtools}
\usepackage{xcolor}
\usepackage[colorlinks=true,linkcolor=blue!55!black,citecolor=blue!55!black,urlcolor=blue!55!black]{hyperref}

\newtheorem{theorem}{Theorem}[section]
\newtheorem{lemma}[theorem]{Lemma}
\theoremstyle{remark}
\newtheorem{remark}[theorem]{Remark}

\title[A022030 Satisfies Its Linear Recurrence]{The Alternating Extremal-Ratio Sequence A022030 Satisfies Its Conjectured Linear Recurrence}
\author{Xiaojian Zeng}
\address{PBC School of Finance, Tsinghua University, 100190}
\email{xiaojianzeng@foxmail.com}
\date{September 2, 2026}
\hypersetup{pdftitle={The Alternating Extremal-Ratio Sequence A022030 Satisfies Its Conjectured Linear Recurrence},pdfauthor={Xiaojian Zeng}}

\begin{document}
\begin{abstract}
Let $a_0=4,a_1=16$.  For even $n$, let $a_{n+2}$ be the greatest integer
whose ratio to $a_{n+1}$ is strictly smaller than $a_{n+1}/a_n$; for odd
$n$, use the least integer with the reverse strict inequality.  We prove
$a_n=4a_{n-1}-a_{n-3}+a_{n-4}$ for every $n\ge4$.  The proof turns the
strict-rounding condition into a signed determinant bound.  A
positive-coefficient recurrence proves the required sign, while a factor-$3$
growth estimate proves the upper remainder bound.
\end{abstract}
\maketitle

\section{Introduction}
The current definition and data of OEIS A022030 are recorded in
\cite{oeisA022030}.  The same alternating rule and recurrence conjecture
appear in the fixed Formal Conjectures source \cite{formalConjectures22030}.
The OEIS history also records an older, nonalternating definition.  A recent
proof summary \cite[Table~1, A022030 row]{adamczewski2026} concerns that older
sequence, not the alternating sequence treated here.

We prove the current conjecture.  The main point is not finite agreement but a
uniform remainder estimate for the candidate linear recurrence.

\section{The rounding criterion}
For positive integers $x,y$, the greatest integer $m$ satisfying $m/y<y/x$
is $\lceil y^2/x\rceil-1$, including when $x$ divides $y^2$.  The least
integer satisfying the reverse strict inequality is $\lfloor y^2/x\rfloor+1$.

Let
\begin{equation}\label{eq:Lrec}
 \begin{gathered}
 L_0=4,\quad L_1=16,\quad L_2=63,\quad L_3=249,\\
 L_n=4L_{n-1}-L_{n-3}+L_{n-4}\quad(n\ge4),
 \end{gathered}
\end{equation}
and set
\begin{equation}\label{eq:defect}
 d_n=L_{n+1}^2-L_nL_{n+2},\qquad e_n=(-1)^n d_n.
\end{equation}

\begin{lemma}\label{lem:rounding}
If $L_n>0$ and $1\le e_n\le L_n$ for all $n\ge0$, then the sequence
defined by the alternating extremal-ratio rule is $L$.
\end{lemma}
\begin{proof}
For even $n$,
\[
 \frac{L_{n+1}^2}{L_n}=L_{n+2}+\frac{e_n}{L_n},
 \qquad 0<\frac{e_n}{L_n}\le1,
\]
hence $\lceil L_{n+1}^2/L_n\rceil-1=L_{n+2}$.  For odd $n$ the sign is
reversed, so $\lfloor L_{n+1}^2/L_n\rfloor+1=L_{n+2}$.  Induction from
$L_0,L_1$ identifies the sequences.
\end{proof}

\section{The defect recurrence}
\begin{lemma}\label{lem:growth}
The sequence $L$ is positive and satisfies $L_n>3L_{n-1}$ for every $n\ge2$.
Consequently $L_{n-j}<L_n/3^j$ whenever $1\le j\le n$.
\end{lemma}
\begin{proof}
The assertion is direct for $n=2,3$.  If the terms through $L_{n-1}$ are
positive and strictly increasing, then, for $n\ge4$,
\[
 L_n-3L_{n-1}=L_{n-1}-L_{n-3}+L_{n-4}>0.
\]
This simultaneously preserves positivity and proves the induction step.  The
last assertion follows by iteration.
\end{proof}

\begin{lemma}\label{lem:defectrec}
The signed defects begin
\[
(e_0,\ldots,e_9)=(4,15,9,105,241,405,2275,4530,12186,46605)
\]
and, for every $n\ge10$, satisfy
\begin{equation}\label{eq:positive-rec}
\begin{split}
231e_n={}&460e_{n-2}+3435e_{n-3}+3585e_{n-5}\\
&+215e_{n-8}+30e_{n-9}+6e_{n-10}.
\end{split}
\end{equation}
In particular, every $e_n$ is a positive integer.
\end{lemma}
\begin{proof}
Direct substitution of \eqref{eq:Lrec} into \eqref{eq:defect} gives
\begin{equation}\label{eq:six-rec}
 e_n=3e_{n-2}+15e_{n-3}-3e_{n-4}+e_{n-6}\qquad(n\ge6).
\end{equation}
Thus
\[
 E(x):=\sum_{n\ge0}e_nx^n
 =\frac{4+15x-3x^2+x^4}{1-3x^2-15x^3+3x^4-x^6}.
\]
The exact polynomial identities
\begin{align*}
&(1-3x^2-15x^3+3x^4-x^6)(231+233x^2+30x^3+6x^4)\\
&\quad=231-460x^2-3435x^3-3585x^5-215x^8-30x^9-6x^{10},\\
&(4+15x-3x^2+x^4)(231+233x^2+30x^3+6x^4)\\
&\quad=924+3465x+239x^2+3615x^3+6x^4+215x^6+30x^7+6x^8
\end{align*}
give \eqref{eq:positive-rec} by coefficient comparison.  Its coefficients and
the first ten values are positive, so strong induction gives $e_n>0$.
\end{proof}

\section{Main theorem}
\begin{theorem}\label{thm:main}
Let $a_0=4,a_1=16$.  For each $n\ge0$, if $n$ is even, let $a_{n+2}$ be the
greatest integer satisfying
\[
 \frac{a_{n+2}}{a_{n+1}}<\frac{a_{n+1}}{a_n};
\]
if $n$ is odd, let $a_{n+2}$ be the least integer satisfying the reverse
strict inequality.  Then all terms are positive and
\[
 a_n=4a_{n-1}-a_{n-3}+a_{n-4}\qquad(n\ge4).
\]
\end{theorem}
\begin{proof}
We prove $e_n\le L_n$.  This holds directly for $0\le n\le9$.  Suppose
$n\ge10$ and the claim holds at smaller indices.  The preceding lemmas give
\[
\begin{split}
231e_n&\le L_n\left(\frac{460}{3^2}+\frac{3435}{3^3}
+\frac{3585}{3^5}+\frac{215}{3^8}+\frac{30}{3^9}+\frac6{3^{10}}\right)\\
&=L_n\frac{3801197}{19683}<231L_n,
\end{split}
\]
where $3801197<231\cdot19683=4546773$.  Strong induction gives
$1\le e_n\le L_n$ for all $n$, because each $e_n$ is a positive integer.
Lemma \ref{lem:rounding} identifies $a=L$, and \eqref{eq:Lrec} proves the
claim.  Positivity follows from Lemma \ref{lem:growth}.
\end{proof}

\begin{remark}
The endpoint $e_n=L_n$ is harmless: it corresponds to an integral quotient
one unit beyond $L_{n+2}$, and the strict ceiling/floor convention still
selects $L_{n+2}$.  Equality occurs at $n=0$ among the base cases; the
induction above is strict for $n\ge10$.
\end{remark}

\section*{Reproducibility and computational disclosure}
The proof is exact and does not depend on finite sampling.  One accompanying
script uses symbolic polynomial expansion to check the universal defect
identity and the displayed generating-function identities.  A separate
standard-library Python program generates the linear and nonlinear sequences
by integer arithmetic, checks the first 2000 terms, and verifies the
certificate recurrence and coefficient comparison.  These computations audit
the algebra and base data; the infinite conclusion is the induction above.
Automated language tools assisted with source discovery, verification-script
drafting, and prose revision.  The author independently checked the cited
sources and all decisive exact identities and is responsible for the
mathematical claims.

\bibliographystyle{amsplain}
\bibliography{references}
\end{document}
